\chapter{Introduction}
\label{chap:intro}

Quantum computing promises to solve certain problems far more efficiently than classical computers. One of the clearest examples is the simulation of quantum systems. The Hilbert-space dimension of a quantum many-body system grows exponentially with the system size, making classical simulation quickly impractical. For example, fully specifying the state of 50 qubits requires $2^{50} \approx 10^{15}$ complex amplitudes \cite{Feynman1982,Lloyd1996}. More broadly, quantum algorithms also offer advantages in areas such as cryptography, optimization, and machine learning \cite{Dalzell_2025}.

Quantum algorithms, such as Shor's algorithm, require fault-tolerant quantum computation~\cite{shor1997polynomial,shor1996faulttolerant}. Quantum error correction enables fault tolerance by encoding a logical qubit into many physical qubits. Logical errors are suppressed by increasing the code size or lowering the physical error rate. In the surface code, error correction only helps when the physical error rate falls below a threshold of order $10^{-2}$~\cite{Fowler2012SurfaceCodes,Campbell2017Roads}. In practice, however, operating below threshold is not sufficient, because the required number of physical qubits still depends strongly on the physical error rate. According to the empirical scaling relation in the surface code~\cite{Fowler2012SurfaceCodes}, for a target logical error rate below $10^{-15}$, reducing the physical error rate from $10^{-3}$ to $10^{-4}$ lowers the required number of physical qubits from about $4.8\times10^{3}$ to about $8.4\times10^{2}$.

Several hardware platforms are being explored for fault-tolerant quantum computation~\cite{Monroe2013,Bruzewicz2019,Kok2007,Slussarenko2019,Doherty2013,Pezzagna2021,Saffman2016,Bluvstein2022}. Superconducting circuits are particularly attractive because they support fast gates, flexible microwave control, and scalable fabrication using mature lithographic methods~\cite{Krantz2019,Kjaergaard2020}. However, superconducting devices are sensitive to noise, which causes decoherence. In addition, the fidelities of quantum operations are degraded by several factors, including imperfect pulses, crosstalk, and leakage. Decoherence and control errors are therefore two major challenges for superconducting qubits. This thesis investigates these challenges.

Coherence in superconducting circuits has been improved dramatically since the first demonstration of coherent control in the Cooper-pair box~\cite{Nakamura1999}. The Cooper-pair box demonstrated basic qubit control, but its strong sensitivity to charge-offset noise limited coherence to the nanosecond scale~\cite{Nakamura2002,Astafiev2004}. The transmon suppresses charge noise by shunting the Cooper-pair box with a large capacitance, thereby greatly improving coherence at the cost of reduced anharmonicity~\cite{Koch2007,Schreier2008,Houck2008}. The fluxonium uses a large inductive shunt, combining protection against charge offsets with strong anharmonicity and high coherence~\cite{Manucharyan2009,Nguyen2019}. More recently, protected circuits such as the $0\text{--}\pi$ qubit, the protomon, and $\cos(2\phi)$-based qubits have sought stronger intrinsic noise protection, albeit with more complex circuit designs~\cite{Brooks2013,Gyenis2021,Kumar2024Protomon,Smith2020}. In parallel, bosonic qubits based on superconducting microwave cavities have emerged as a complementary approach in which quantum information is stored in long-lived linear oscillator modes~\cite{Cai2021,Copetudo2024,Paik2011,Reagor2013,Reagor2016,Milul2023,Romanenko2020,Joshi2021}. These systems are attractive because bosonic encodings can reduce the hardware overhead required for quantum error correction. However, preserving coherence in cavity-based architectures remains challenging because bosonic control typically relies on nonlinear elements that introduce additional decoherence channels. Understanding and mitigating these inherited decoherence channels is essential for realizing the potential of cavity-based architectures.

Control methods have also evolved. Early demonstrations of qubit control often relied on simple analytic pulse optimization. As superconducting processors became more complex, however, this approach became less robust. Weak anharmonicity, spectral crowding, and crosstalk limit the reliability of analytic methods~\cite{Motzoi2009DRAG,Dumas2024}. This has motivated numerical optimal-control methods that can incorporate these constraints directly~\cite{Brif2010Control,Khaneja2005,deFouquieres2011GRAPE}. One widely used method is Gradient Ascent Pulse Engineering (GRAPE). With automatic differentiation, GRAPE can evaluate cost-function gradients automatically during optimization. However, it can require substantial memory, which limits its application to large quantum systems.

This thesis addresses two challenges in superconducting quantum hardware: efficient optimal control and coherence engineering. Chapter~\ref{chap:sc-qubits} provides background on superconducting qubits, decoherence, and control. Chapter~\ref{chap:grape} presents the optimal-control results. It studies gradient-based pulse optimization and compares three methods for evaluating cost-function gradients in terms of memory use and runtime. It also provides practical guidance for selecting an efficient gradient-evaluation method for a given task. Chapter~\ref{chap:stark-sweet-spots} presents the coherence-engineering results for cavity-based superconducting architectures. High-$Q$ cavities can exhibit negligible intrinsic dephasing. However, their control often depends on nonlinear Josephson elements. Frequency fluctuations in these elements are inherited by the cavity and cause dephasing. We show that a weak off-resonant drive applied to the nonlinear element creates an AC Stark shift, leading to a Stark-induced sweet spot. This mechanism extends the cavity dephasing time by more than an order of magnitude. Finally, Chapter~\ref{chap:conclusion} summarizes the main results and discusses future directions.